% Phase 2 Research Proposal  (course limit: at most TWO pages, excluding references)
% Compile: pdflatex proposal.tex  (twice, for references)
% All quantitative claims are drawn from experiments/*.json and are audit-verified.
\documentclass[10pt]{article}

\usepackage[
    top=0.3in,
    bottom=0.8in,
    left=0.75in,
    right=0.75in
]{geometry}
\usepackage{times}
\usepackage{amsmath,amssymb}
\usepackage{hyperref}
\usepackage{enumitem}
\usepackage{booktabs}
\usepackage{xcolor}

\setlist{nosep,leftmargin=*}

\title{\textbf{Research Proposal}\\
Causal Reinforcement Learning with Hidden Confounders and Partial Observability}
\author{
Team Members: Soheil Sayah Varg, Mahdi Shirinbayan, Mahdi Ostadmohammadi
}
\date{}

\begin{document}
\maketitle

% =====================================================================
\section{Motivation}

Offline reinforcement learning from logged data is unreliable when the behavior policy acted on information the
learner never sees. In a partially observable MDP (POMDP) whose \emph{latent} state drives both the logged
actions and the outcomes, the data are \emph{confounded}: a learner that treats the observation as the state
systematically mis-estimates policy value, and no amount of data fixes the bias. This is the central obstacle
to trustworthy offline RL in healthcare, operations, and any domain with unrecorded drivers of decisions.

Hong, Qi \& Xu (ICML 2024)~\cite{hong2024} give an elegant theory for exactly this setting: using a
pre-decision \emph{negative control} $O_0$, they identify the confounded POMDP's reward and transition models
through two families of \emph{bridge functions} ($b_R$, $b_D$), then plan pessimistically over a
policy-independent confidence region. The paper is, however, \emph{theory only} --- no algorithm, no code, no
experiments. Its neighbors are likewise unrealized empirically for this exact pipeline: the closest follow-up
runs the bridge \emph{estimator} in isolation~\cite{venkatesh2026}, and the pessimism precedent is
theory-only and for the \emph{un}confounded case~\cite{guo2022}. Whether this appealing identification theory
actually produces usable value estimates and policies at practical sample sizes is therefore open --- and is
what we set out to test.

% =====================================================================
\section{Research Question}

\emph{When the Hong--Qi--Xu confounded-POMDP framework is implemented end-to-end, does model-based bridge
identification actually de-bias policy evaluation and yield good policies at practical sample sizes, and can
its abstract pessimistic $\max$--$\min$ step be made computationally tractable?} Because the answer turns out
to be \textbf{regime-dependent}, our concrete deliverable is to \emph{characterize where} de-biasing helps
(strong confounding, exact identification) versus where estimation variance dominates and the method is beaten
by a naive baseline (benchmark scale).

% =====================================================================
\section{Proposed Direction}

\noindent\textbf{Contribution (capability, not primacy).} We build what is, to our knowledge, the first
end-to-end empirical implementation that jointly estimates \emph{both} bridge families ($b_R$, $b_D$) via the
paper's two-stage procedure and carries them through plug-in value identification to a tractable pessimistic
policy-\emph{selection} step. We do \emph{not} claim primacy over concurrent model-based confounded-POMDP work
(e.g.\ DPC-MBRL, which we have not reviewed). Per block, the pessimistic inner minimization is the standard
ellipsoid support-function identity (exact because the stage-2 risk is quadratic and the plug-in value
multilinear); our genuinely new content is a \emph{diagnosis} of a failure mode in the coupled multi-block
minimization and an honest account of a signal-subspace remedy.

\begin{itemize}
  \item \textbf{Estimation (faithful).} Stage~1: conditional mean embedding $\hat\mu_{W_t|x}$ by kernel ridge
  regression $+$ conditional density/pmf targets $\hat p(y_t|x_t)$. Stage~2: closed-form bridge solve
  $\alpha = (G + N_2\lambda_2 I)^{-1}\hat p$, $G = (\Gamma^\top K_W \Gamma) \odot K_{Y''}$. \emph{Methods
  finding:} the paper's RKHS hyperparameter schedules do \emph{not} transfer to the tabular class; we
  grid-calibrate split $0.7{:}0.3$ ($\approx 2.3{:}1$, consistent with $N_1 \ge N_2$), $\lambda_1 = 1/N_1$,
  $\lambda_2 = 0.03/\sqrt{N_2}$. Ill-posedness is monitored by the restricted-spectrum signal eigenvalue
  $\hat\sigma_2$, \emph{not} $\mathrm{cond}(G)$ (which the structural null space pins flat).
  \item \textbf{Tractable pessimism.} Confidence regions are exact ellipsoids centered at $\hat b$; the
  per-block $\min$ is $V(\hat b) - \sqrt{\xi}\,\|g_\pi\|_{H^{-1}}$. The hard part --- the \emph{coupled}
  multilinear $\min$ across blocks --- is solved by coordinate descent (reward blocks global, dynamic blocks
  heuristic), so $V_{\mathrm{low}} \le V_{\mathrm{true}}$ is verified empirically, not certified.
  \item \textbf{Environments/tools.} \emph{Primary:} the public \texttt{clinicalml/gumbel-max-scm} tabular
  benchmark~\cite{oberst2019} (720 observable-core $\times$ 2 latent states, 8 actions, $T{=}3$; the latent
  index drives dynamics/behavior but is censored from logs; a manufactured pre-decision $O_0$ satisfies the
  negative-control assumption by construction, and a per-step noisy marker guarantees bridge existence ---
  deterministic masking gives rank $720<1440$, no one-step bridges). \emph{Companion:} a minimal toy POMDP
  ($|S|{=}2,|A|{=}2,|O|{=}3,T{=}3$) with exact oracle bridges --- the only substrate where recovery error is
  measurable. Ground truth by dynamic programming on the latent SCM; NumPy, no GPU.
\end{itemize}

% =====================================================================
\section{Proof of Concept}

The pipeline is \textbf{already implemented and oracle-verified end-to-end}; all numbers below are from
\texttt{experiments/*.json}.

\begin{itemize}
  \item \textbf{Oracle verification.} On the toy, true bridges reproduce the Theorem-3.5 chain value vs exact
  DP to $4.4{\times}10^{-16}$; the population-limit solve recovers the oracle bridges to $3.0{\times}10^{-15}$;
  primal/dual closed forms agree to $2.4{\times}10^{-15}$.
  \item \textbf{Bridge recovery.} L2 error vs oracle decreases monotonically, $\mathrm{err}_{b_R}: 1.36\!\to\!
  0.46$ and $\mathrm{err}_{b_D}: 0.73\!\to\!0.22$ over $N{=}1\mathrm{k}\!\to\!20\mathrm{k}$ (log-log OLS slopes
  $N^{-0.36}$, $N^{-0.39}$; 3 seeds, per-seed spread persisted, CIs not yet reported).
  \item \textbf{De-biasing (toy, where it works).} The plug-in beats the confounding-blind naive evaluator in
  \textbf{12/12} $(N,\text{policy})$ cells; worst-policy plug-in bias $0.11\!\to\!0.03$
  ($N{=}5\mathrm{k}\!\to\!50\mathrm{k}$) while naive stays $\approx 0.23$. (Naive bias itself is
  $+0.14$ to $+0.33$ per policy, $8.1$--$22.0\%$ relative, flat in $N$ --- the target signal.)
  \item \textbf{Pessimism --- the honest headline.} \emph{(Diagnosis, verified)} the literal $\max$--$\min$ is
  uninformative at every width that \emph{contains} the truth: where the full-space region covers the true
  bridges ($c\ge 0.1$), $V_{\mathrm{low}}$ blows up ($-1.12$ at onset, to $-1755$; true values
  $\approx 1.5$--$2.2$) and mis-selects (regret $0.68$), because $|O|{>}|S|$ null directions carry ellipsoid
  weight only $\lambda_2$. \emph{(Remedy is a heuristic, not a fix)} restricting to the identified signal
  subspace keeps $V_{\mathrm{low}}$ informative ($1.35$ at $c{=}0.03$) and selects correctly (regret $0$, all
  seeds), \emph{but} the true bridges leak $\approx 21\%$ of their norm out of the subspace, so the projected
  region never contains the truth (honest joint coverage $=0$): a genuine \emph{coverage/informativeness
  tension}, which we take as the honest finding rather than an overclaimed ``repair''.
  \item \textbf{Benchmark boundary (negative, disclosed).} At $720{\times}2$ scale ($N\le 60\mathrm{k}$) the
  plug-in is \emph{worse} than naive in \textbf{15/16} cells (14/16 vs the observable empirical naive); the harm
  is largest under \emph{zero} confounding ($23$--$46\times$ the naive bias), where there is nothing to correct
  --- so per-core estimation \emph{variance}, not the confounding signal, dominates at feasible $N$.
\end{itemize}

% =====================================================================
\section{Risks and Fallback Plan}

\begin{itemize}
  \item \textbf{What might fail / is uncertain.} (i) \emph{Pessimism guarantees:} the coupled inner $\min$ is
  heuristic and the signal projection excludes part of the truth, so $V_{\mathrm{low}}\le V_{\mathrm{true}}$ is
  empirical, not certified --- we report leak, restart gaps, and the tension rather than overclaim. (ii)
  \emph{Estimation variance at scale:} the benchmark negative result shows the estimator dominated by naive at
  feasible $N$; the identified next step is cross-fitting / partial pooling to reduce the ratio/inversion
  variance of the per-cell solve. (iii) \emph{Hyperparameters} are not computable from theory --- we
  cross-validate $\lambda_1,\lambda_2$ and sweep $\xi$ against coverage curves, never claiming a theoretical
  $\xi$. (iv) \emph{Baselines:} only a confounding-blind naive baseline is in place; adding a confounding-aware
  proximal-OPE comparator is a rigor goal.
  \item \textbf{Fallback ladder.} tabular-exact only $\to$ pessimistic \emph{selection} over fixed candidates
  (avoids the intractable full optimization) $\to$ \emph{floor deliverable}: the oracle-verified bridge
  estimation and toy de-biasing curves, which stand on their own as the first empirical characterization of
  this framework, plus the diagnostic pessimism findings above.
\end{itemize}

% =====================================================================
\begin{thebibliography}{9}
\small
\bibitem{hong2024} M.~Hong, Z.~Qi, Y.~Xu.
\textit{Model-based Reinforcement Learning for Confounded POMDPs}. ICML, 2024.
\bibitem{venkatesh2026} V.~Venkatesh, A.~Malikopoulos.
\textit{Cross-fitted Proximal Learning for Model-Based Reinforcement Learning}. arXiv:2604.05185, 2026.
\bibitem{guo2022} X.~Guo, M.~Cai, Z.~Yang, Z.~Wang.
\textit{Provably Efficient Offline Reinforcement Learning for Partially Observable Markov Decision Processes}.
ICML, 2022.
\bibitem{oberst2019} M.~Oberst, D.~Sontag.
\textit{Counterfactual Off-Policy Evaluation with Gumbel-Max Structural Causal Models}. ICML, 2019.
\bibitem{shi2022} C.~Shi, M.~Uehara, J.~Huang, N.~Jiang.
\textit{A Minimax Learning Approach to Off-Policy Evaluation in Confounded POMDPs}. ICML, 2022.
\bibitem{bennett2024} A.~Bennett, N.~Kallus.
\textit{Proximal Reinforcement Learning: Efficient Off-Policy Evaluation in POMDPs}.
Operations Research, 2024.
\bibitem{mastouri2021} A.~Mastouri et al.
\textit{Proximal Causal Learning with Kernels}. ICML, 2021.
\end{thebibliography}

\end{document}
